%%%%%%%%%%%%%%%%%%%%%%%%%%%%%%%%%%%%%%%%%%%%%%
%% 02_testing_a_signal.tex --- 正文层 (BODY LAYER)
%%
%% 这里只有内容。字体、边距、颜色、定理环境长什么样，全在
%% ../../preamble.tex 里；改版式不用碰这个文件。
%%
%% 这是正文片段，没有 \documentclass 也没有 document 环境：
%% ../../build.sh 单独编译它，all_chapters.tex 把它 \input 进合订本。
%%%%%%%%%%%%%%%%%%%%%%%%%%%%%%%%%%%%%%%%%%%%%%

\chapterhead{02}{Testing a Signal}{02_testing_a_signal}


\section{A signal is a hypothesis, not a formula}\label{sec:signal-hypothesis}

\subsection{The three objects}\label{sec:three-objects}

Everything in this chapter is built from three numbers, one of each per asset per date.

\begin{tabularx}{\linewidth}{|>{\raggedright\arraybackslash}p{0.16\linewidth}|>{\raggedright\arraybackslash}p{0.10\linewidth}|X|>{\raggedright\arraybackslash}p{0.17\linewidth}|}
\hline
\textbf{Object} & \textbf{Symbol} & \textbf{What it is} & \textbf{Known at $t$?} \\ \hline
\textbf{Weight}
  & $w_{s,t}$
  & the share of capital held in asset $s$ on date $t$, signed---negative is a short
  & \textbf{yes}, you choose it \\ \hline
\textbf{Signal}
  & $\mathit{MOM}_{s,t}$
  & a number computed from data available at $t$, meant to say something about what comes next
  & \textbf{yes}, you compute it \\ \hline
\textbf{Forward return}
  & $r_{s,t}$
  & what the asset then goes on to deliver while the position is held
  & \textbf{no}---this is the unknown \\ \hline
\end{tabularx}

\begin{note}[Scope]
This chapter describes \textbf{one asset at a time}. That is not a simplification: assets have
different volatilities, so their raw signals are not on one scale and cannot be compared until
\secref{sec:risk-adjusted} puts them there---and combining several into one book comes after that,
in the position-construction step.
\end{note}

\begin{definition}[Binary momentum]\label{def:binary-momentum}
The simplest member of the family---the sign of last period's return, and nothing else:
\[
  \mathit{MOM}_{s,t} = \operatorname{sign}\left( r_{s,t-1} \right)
\]
where $s$ indexes the asset and $t$ the date, $r_{s,t-1}$ is that asset's return over the period
just ended, and $\mathit{MOM}_{s,t}$ is either $+1$ (long) or $-1$ (short), with nothing in between.
\end{definition}

\textbf{Trend direction survives; trend strength is thrown away.} A window in which the asset rose
20 percent and one in which it rose 10 percent produce the \emph{same} signal, so a strategy built
on it takes the same size in both.

\fig{binary-momentum}{Left, one asset's price over two lookback windows, rebased to 100: one climbs
to 120 and the other to 110. Right, the signal each path produces---two bars of identical height at
$+1$. The magnitude of the move does not reach the signal. Illustrative.}{fig:binary-momentum}

\begin{definition}[Momentum]\label{def:momentum-avg}
Averaging over $N$ periods rather than one, and keeping the value rather than its sign:
\[
  \mathit{MOM}_{s,t} = \operatorname{Avg}\left( r_{s,t-i} \right)
    = \frac{1}{N} \sum_{i=1}^{N} r_{s,t-i}, \qquad i = 1 \ldots N
\]
with $N$ the lookback length and $i$ the lag inside it---the average runs over the $N$ returns
ending yesterday, so at $N = 21$ it is last month's average daily move.
\end{definition}

\begin{example}\label{ex:three-assets}
Three assets over a five-day lookback:

\begin{center}
\begin{tabular}{|c|l|r|r|}
\hline
\textbf{Asset} & \textbf{Last five returns} & \textbf{Avg $\rightarrow$ the signal} & \textbf{After $\operatorname{sign}$} \\ \hline
A & $+3\%$, $+2\%$, $+4\%$, $+1\%$, $+2\%$      & $\mathbf{+0.024}$ & $+1$ \\ \hline
B & $+1\%$, $-0.5\%$, $+1\%$, $0\%$, $+0.5\%$   & $\mathbf{+0.004}$ & $+1$ \\ \hline
C & $-2\%$, $-3\%$, $-1\%$, $-2\%$, $-2\%$      & $\mathbf{-0.020}$ & $-1$ \\ \hline
\end{tabular}
\end{center}

The right column cannot tell A from B; the left makes A \textbf{six times} B, and since
\secref{sec:signal-is-portfolio} sets each weight proportional to its signal, A is held six times
larger. That ratio is the magnitude, and it is all the dropped $\operatorname{sign}$ was costing.
\end{example}

\begin{note}[The level is not yet meaningful]
$0.024$ means nothing on its own---2.4 percent in a calm year and 2.4 percent in a violent one are
not the same trend. Only \emph{relative} sizes are ever used, and \secref{sec:risk-adjusted} is what
puts them on one scale.
\end{note}

\subsection{Why studying the signal is the same as studying the portfolio}\label{sec:signal-is-portfolio}

A signal is not a strategy: it earns nothing and cannot be held. The reason to spend a chapter on
one runs backwards, from what you are actually paid on---\textbf{so start from the portfolio,
because that is what you optimize.}
\[
  R_t = \sum_s w_{s,t} r_{s,t}
\]
$R_t$ is the portfolio's return on date $t$, summed over every asset in the universe. Of the two
factors in each term only one is yours---$r_{s,t}$ is the market's answer, and it arrives after the
weight is already set. \textbf{So the whole of portfolio construction is the choice of $w_{s,t}$.}

\textbf{But the weights are not free.} Two sums over the book are fixed before any signal is
consulted---the \textbf{net}, fixing how much \emph{market} the book carries since longs and shorts
cancel in it, and the \textbf{gross}, fixing how much \emph{leverage} is deployed since they add:
\[
  \textbf{net:} \quad \sum_s w_{s,t} = 1
  \qquad\qquad
  \textbf{gross:} \quad \sum_s \left| w_{s,t} \right| = G
\]
with $G$ the gross target---200 percent on a 150/50 book---and the net holding at 1 only to within a
tolerance $\delta$, since rebalancing is discrete. The two are independent: one asset at 100 percent
and a book long 200 percent / short 100 percent carry the same net and three times the position,
which is why gross is never free (\chapref{1}{01_what_is_cta}).

\textbf{The weights are where the signal enters.} You want weight where the forward return is about
to be high, and that is precisely the number you do not have. The signal is the stand-in you put in
its place, which means setting $w_{s,t} \propto \mathit{MOM}_{s,t}$.

\begin{claim}\label{clm:rho-is-everything}
Under that rule the portfolio's expected return is proportional to the correlation between the
signal and the forward return, and to nothing else that can change its sign.
\end{claim}

\begin{proof}
Take the signal and the return as centred. With $w_{s,t} \propto \mathit{MOM}_{s,t}$,
\[
  E[R_t] \propto \sum_s E\left[ \mathit{MOM}_{s,t}\, r_{s,t} \right]
    = \sum_s \rho_s\, \sigma_{\mathit{MOM},s}\, \sigma_{r,s}
\]
using $E[xy] = \rho\, \sigma_x \sigma_y$ for centred $x$ and $y$, where $\rho_s$ is the correlation
between asset $s$'s signal and its forward return and $\sigma_{\mathit{MOM},s}$, $\sigma_{r,s}$
their standard deviations. The dropped constant is a leverage choice and the volatilities are
measurable properties of the asset---all three positive whatever the signal does. \textbf{Only
$\rho_s$ can carry a sign, so only $\rho_s$ decides whether the portfolio makes money or loses it.}
\end{proof}

\begin{note}[The constraints do not disturb this]
Reaching a legal weight vector means \textbf{shifting} the signal until the net lands right and
\textbf{scaling} until the gross does---both affine with a positive scale, under which correlation
is unchanged: $\rho(a + bx, y) = \rho(x, y)$ for $b > 0$. The $\rho$ measured on the raw signal is
therefore the $\rho$ governing the constrained portfolio, which is what lets the rest of the chapter
ignore weights altogether; the position-construction step carries out both operations.
\end{note}

Everything collapses to that one number, and it is a number about the signal alone:
\[
  \textbf{hypothesis:} \quad \mathit{MOM} \uparrow \Longrightarrow \text{return} \uparrow
  \qquad\qquad
  \textbf{that is:} \quad \rho > 0
\]

\textbf{Write that down before the code}---it names what would falsify the signal, and one you
cannot falsify is a plot you will rationalize either way. Two things follow from it, and between
them they set the shape of the rest of the chapter:

\begin{tabularx}{\linewidth}{|X|>{\raggedright\arraybackslash}p{0.22\linewidth}|}
\hline
\textbf{Consequence} & \textbf{Dealt with in} \\ \hline
\textbf{The portfolio never has to be built to test the signal.} Measure $\rho$ between the signal
and the forward return directly, and you have measured the strategy
  & \secref{sec:hoped-vs-got}, \secref{sec:two-ways-out} \\ \hline
\textbf{At equal $\rho$ a volatile asset contributes more than a quiet one}, since each term carries
$\sigma_{\mathit{MOM},s} \sigma_{r,s}$ alongside $\rho_s$. Ranking raw signals therefore ranks
volatilities
  & \secref{sec:risk-adjusted} \\ \hline
\end{tabularx}


\section{What you hoped for, and what you get}\label{sec:hoped-vs-got}

\subsection{Which return the signal is scored on}\label{sec:which-return}

Before signal and return can be plotted against each other, \emph{the forward return} has to be
pinned down: \secref{sec:signal-is-portfolio} measures $\rho$ against it, but the phrase is not yet
precise enough to compute. Three spans sit on the timeline, in this order:

\begin{tabularx}{\linewidth}{|>{\raggedright\arraybackslash}p{0.18\linewidth}|X|>{\raggedright\arraybackslash}p{0.22\linewidth}|}
\hline
\textbf{Span} & \textbf{What it holds} & \textbf{Typical size} \\ \hline
\textbf{Lookback}
  & the $N$ returns the signal averages, ending at $t-1$
  & 20--250 periods \\ \hline
\textbf{Gap}
  & $g$ periods thrown away---neither averaged into the signal nor scored against it
  & 0, a week, or a month \\ \hline
\textbf{Paired return}
  & $r_{s,t+g}$, the one period the signal is judged on
  & one period \\ \hline
\end{tabularx}

\fig{signal-return-alignment}{Above, one observation laid on a timeline of one cell per period: the
lookback the signal averages, then the discarded gap $g$, then the single period the signal is
scored on. Below, the same pattern for dates $t$, $t+1$ and $t+2$, each shifted one cell right---the
lookbacks overlap in all but one cell, while the three scored cells form a descending diagonal, so
no two dates share a scored return. Illustrative.}{fig:signal-return-alignment}

The gap is the only free choice of the three---the $-1$ in \textbf{12-1 momentum} is exactly this
gap, a twelve-month window stopping one month short of today.

\textbf{Why leave one at all.} Not executability: \secref{sec:three-objects}'s window already ends
at $t-1$, so $g = 0$ is knowable in time and anything tighter is look-ahead
(\secref{sec:info-availability}). Executability fixes the floor at zero; \textbf{reversal} decides
everything above it. A trend that has just formed hands a little of it back, and at daily frequency
that rebate is large enough to cancel the edge outright---which is how a real signal arrives at the
bar plot looking flat, or upside down. The mechanisms behind it are this chapter's Background.

\textbf{Both directions cost something}, which is what makes $g$ a \textbf{hyperparameter} rather
than a constant. Reversal strengthens as the sampling gets finer---the bid--ask bounce is a roughly
fixed number of ticks per trade while the trend grows with the horizon---but a wider gap also throws
away the opening of the trend. Set it from your trading frequency and a prior on how long the rebate
lasts; running it at a day, a week and a month and keeping whichever bar plot looks best is what
\chapref{5 \midot{} Modelling}{05_modelling} warns about.

\begin{note}[Neighbouring observations are near-copies]
Step one date forward and the lookback keeps all but one of its cells, as the figure's lower panel
shows. Consecutive signals are therefore not independent draws, which is where
\secref{sec:bar-plot-trend}'s caveat on the effective sample size comes from.
\end{note}

\subsection{The plot everyone draws first}\label{sec:first-plot}

\secref{sec:signal-hypothesis}'s hypothesis says forward return rises with the signal, so the first
move is to plot one against the other. The picture you had in mind is the leftmost panel below. What
comes back---for well over 99 percent of the scatters you will ever draw---is the rightmost.

\fig{scatter-ladder}{Four scatters of forward return against a signal, both axes in standard
deviations. The first three are drawn at correlations of 80, 45 and 30 percent from synthetic
points; by 30 percent the tilt is barely arguable, which is the eye's floor. The fourth is measured
on \texttt{CTA\_data} at $-1.9$ percent---a formless cloud with a denser core, being real returns
with fatter tails.}{fig:scatter-ladder}

\subsection{Why the real one is a cloud}\label{sec:why-cloud}

\begin{claim}\label{clm:scatter-useless}
The scatter can neither confirm nor refute a signal, because the correlation a working signal
carries sits far below the threshold at which the eye resolves a trend.
\end{claim}

\begin{proof}
Both numbers are readable off the panels above: the eye stops resolving a tilt somewhere around
30 percent, and the fourth panel---21-day risk-adjusted momentum against the next session's return,
measured over the 37 ETFs of \chapref{100 \midot{} The Dataset}{100_dataset}---comes in at
$\mathbf{-1.9}$ \textbf{percent}. Since $1.9\% < 30\%$, a signal that works and a signal that does
not produce the same picture---\textbf{the absence of a visible trend is not evidence of anything.}
\end{proof}

So why is the correlation so small? Because the quantity being plotted is mostly not the signal.

\begin{claim}\label{clm:r-squared}
Split the forward return into the part the signal reaches and the part it does not,
$y = \beta x + \epsilon$. If $\epsilon$ is uncorrelated with $x$, the signal owns exactly $\rho^2$
of the return's variance, where $\rho$ is their correlation.
\end{claim}

\begin{proof}
Uncorrelated components add their variances,
\[
  \operatorname{Var}(y) = \beta^2 \operatorname{Var}(x) + \operatorname{Var}(\epsilon)
\]
and on one regressor with an intercept the explained share is $R^2 = \rho^2$. At $\rho = -0.019$,
$R^2 = 0.00035$---which does not mean ``right 0.035 percent of the time'', but: of the variation in
forward return from one observation to the next, \textbf{0.035 percent is linear in the signal and
99.965 percent is not.}
\end{proof}

\begin{example}\label{ex:bp-split}
Everything below is \textbf{one asset-date's} forward return, never a portfolio's, and every number
is measured over the 57,649 asset-dates of the sample. The daily return standard deviation is
$\sigma_y = 0.0149$---149 bp, a \textbf{basis point} being 0.01 percent:

\begin{center}
\begin{tabular}{|l|c|r|}
\hline
\textbf{Component} & \textbf{Size} & \textbf{At $\rho = -1.9\%$} \\ \hline
Total return          & $\sigma_y$                     & 149 bp \\ \hline
What the signal moves & $\rho \sigma_y$                & \textbf{2.8 bp} \\ \hline
What it does not      & $\sigma_y (1 - \rho^2)^{1/2}$  & \textbf{149 bp} \\ \hline
\end{tabular}
\end{center}

The third row is not a rounding slip: at this correlation the signal removes so little that the
unreachable part is the whole of it to three figures. Across an $x$-axis running from $-2\sigma_x$
to $+2\sigma_x$ the trend line moves about 11 bp end to end, while at every $x$ the points scatter
over roughly 596 bp. \textbf{A 0.11 percent slope inside a 6 percent cloud}---that ratio \emph{is}
the picture.
\end{example}

\begin{note}[Where a larger correlation would come from, and where it would not]
A predictor correlating 50 percent with next period's return would be arbitraged away long before
you found it in anything liquid, so single digits is a competitive market's ceiling rather than a
shortfall in craft. Two things do legitimately raise it and neither is on offer here: a
\textbf{longer horizon}, since the signal's edge accumulates while the noise grows only as the
square root of time, and \textbf{combining many signals}, since a composite reaches correlations no
ingredient has on its own. Quoted numbers near 10 percent are almost always one of those two. A
single signal against a single session is the hardest case in the family, and this chapter is
deliberately standing in it.
\end{note}


\section{Two ways out of the cloud}\label{sec:two-ways-out}

The scatter is unreadable, but the signal may still be in it. There are two things to try, and only
the second one works.

\subsection{Turn down \texttt{alpha}}\label{sec:alpha-opacity}

\texttt{alpha} is a marker's opacity. Turn it far enough down---\texttt{alpha=0.05} on a few
thousand points, lower still on more---and a single point becomes nearly invisible, so only
\emph{overlap} renders and the chart is a density map rather than a mass of ink.

\begin{note}[Not the other alphas]
matplotlib's opacity keyword: not a regression intercept, not the excess return a manager is paid
for, not \chapref{3}{03_shaping_the_lookback}'s smoothing constant $\alpha$. Code font is the tell.
\end{note}

\fig{alpha-opacity}{The same measured asset-dates twice, both axes standardized. Left, at full
opacity, they render as one solid disc with no internal structure. Right, the identical points at an
opacity of 0.01: a smooth round density, denser in the core, with no thin region, no empty corner
and no visible tilt anywhere.}{fig:alpha-opacity}

Sometimes it is enough. A corner bitten out of the cloud---high signal never followed by a bad
loss---would be a finding, because a signal that only rules something \emph{out} is still tradeable.
This data has no such corner, and that is the ordinary outcome.

\textbf{Why it usually fails.} Two reasons that compound:

\begin{itemize}
  \item \textbf{The density comes back smooth.} Tens of thousands of asset-date cells average into a
        clean bivariate blob, and opacity renders that faithfully---a smooth density has no feature,
        and \secref{sec:hoped-vs-got}'s band is fifty times taller than the tilt hiding in it.
  \item \textbf{It treats the symptom.} The scatter spends its resolution on individual noisy points
        when the claim is about their \textbf{average}; no rendering choice changes what is being
        rendered.
\end{itemize}

\begin{note}[Look first anyway]
Always draw it. On the rare occasion something \emph{is} readable---a curve, an empty corner---it
beats every statistic, because it gives the \emph{shape}. The mistake is concluding anything
\textbf{from} a cloud.
\end{note}

\subsection{Beta Method}\label{sec:beta-method}

Stop asking each point what it did, and ask each \emph{group} of points what they did on average.
The answer comes back as a \textbf{bar plot}, and the rest of \secref{sec:two-ways-out} is about
reading that one chart.

\subsubsection{What a bar plot is}\label{sec:bar-plot}

\begin{definition}[Bar plot]\label{def:bar-plot}
Five bars. The horizontal axis is the \textbf{signal}, coarsened into five ordered slots
G1 \ldots G5; the vertical axis is the \textbf{mean forward return} of the observations filed into
each slot. Every observation lands in exactly one bar, and the error bar on it is the standard error
of that mean.
\end{definition}

\begin{note}[A bucket holds dates, not assets]
G5 is not a set of assets. It holds the \textbf{dates} on which this one asset's signal stood in the
top fifth of its own history, and G1 the dates on which it stood in the bottom fifth. One asset
passes through every bucket as the years go by, and the whole chart describes that asset alone.
\end{note}

Five steps build it:

\begin{enumerate}
  \item \textbf{Fix one asset.} Its history gives one signal $\mathit{MOM}_{s,t}$ and one forward
        return $r_{s,t}$ per date.
  \item \textbf{Rank the dates by signal value}, highest to lowest.
  \item \textbf{File each date} into one of five slots by that rank: the highest fifth into G5, down
        to G1.
  \item \textbf{Record} under each slot the forward return that date went on to deliver.
  \item \textbf{Average} the returns filed under each slot.
\end{enumerate}

\begin{note}[Step 2 is where this chapter's one defect lives]
Ranked on \emph{what}, exactly? Take the raw momentum value and everything above still runs---but
the asset's own volatility changes over the years, so $+2$ percent is a strong month in a calm tape
and a quiet one in a violent tape, and step 3 files them into the same slot as if they were the same
event. \secref{sec:risk-adjusted} is the correction, and until it lands the plot below should be
read as machinery rather than as a result.
\end{note}

\fig{bucket-construction}{Six panels; the upper row is illustrative and the lower is measured on
\texttt{CTA\_data}. Left, the population of dates with fifteen drawn at random marked: on a perfect
relationship they sit along a line, on the measured cloud anywhere. Middle, six draws of five dates
plotted against the slots G1 to G5: on the line every draw rises monotonically, on the cloud the six
tangle. Right, the answer with $\pm 2$ standard errors---a clean rising staircase above, and below
the real chart, where G1 stands well clear of its error bar and G2 to G5 sit within a few basis
points of zero, so the staircase descends.}{fig:bucket-construction}

\begin{note}[What makes 57,649 dates from 37 assets one pool]
The measured row does not rank raw momentum across assets. Every signal entering it is already
divided by that asset's own trailing volatility and ranked against that asset's own history before
it lands in a bucket---the route \secref{sec:risk-adjusted} works out in full, with a worked
two-asset example. G1 \ldots G5 are therefore \textbf{percentile} slots, comparable across a
Treasury ETF and a semiconductor ETF for the same reason they are comparable across 2021 and 2023:
nothing but rank survives the trip, so pooling assets costs no more than pooling dates already did.
\end{note}

\textbf{The sort key is the signal, never the return.} Step 2 orders by $\mathit{MOM}_{s,t}$, which
is known at $t$; step 3 only \emph{records} what followed. G1 therefore holds the
lowest-\textbf{signal} observations, not the worst performers.

\begin{example}\label{ex:sort-key}
Five dates from one asset's history, invented, filed both ways:

\begin{center}
\begin{tabular}{|c|r|r|c|c|}
\hline
\textbf{Date} & \textbf{Signal} & \textbf{Forward return} & \textbf{Slot by signal} & \textbf{Slot by return} \\ \hline
$t_1$ & $+2.1\%$ & $-0.4\%$ & G4 & G2 \\ \hline
$t_2$ & $-1.8\%$ & $+0.9\%$ & G1 & G4 \\ \hline
$t_3$ & $+0.6\%$ & $+1.3\%$ & G3 & G5 \\ \hline
$t_4$ & $+3.4\%$ & $+0.2\%$ & G5 & G3 \\ \hline
$t_5$ & $-0.9\%$ & $-1.1\%$ & G2 & G1 \\ \hline
\end{tabular}
\end{center}

Read down \emph{slot by signal} and the returns come out unordered---which is exactly what five
dates look like at a correlation this small. Read down \emph{slot by return} and the ordering is
perfect, and would be perfect for \textbf{any} signal including one straight out of a random number
generator, because each bar is then reporting the sort key back to you. \textbf{The staircase is
evidence only because the thing sorted on and the thing measured are different, and the second was
not knowable when the first was computed.}
\end{example}

One bar therefore carries three pieces of information: which fifth of the signal's range it stands
for, the average return that fifth went on to earn, and how much of that average is sampling noise.

\subsubsection{How the bar plot shows a trend}\label{sec:bar-plot-trend}

\textbf{The trend is the ordering, not the heights.} \secref{sec:signal-hypothesis}'s hypothesis was
\emph{higher signal, higher forward return}; on a bar plot that is precisely the statement
G1 $<$ G2 $<$ G3 $<$ G4 $<$ G5. A single bar's level moves with where the boundaries are cut
(\secref{sec:risk-adjusted})---the ordering is what does not.

Two further readings carry information, and nothing else on the chart does:

\begin{tabularx}{\linewidth}{|>{\raggedright\arraybackslash}p{0.28\linewidth}|X|}
\hline
\textbf{Read} & \textbf{What it establishes} \\ \hline
\textbf{G5 $-$ G1, measured against the error bars}
  & the size of the edge next to what noise alone would draw. A rise that fits inside one error bar
    is not evidence of anything, and \secref{sec:staircase-real} turns that comparison into a
    test \\ \hline
\textbf{The direction of the slope}
  & a \emph{descending} staircase is not a dead signal---it is the same edge carrying the opposite
    sign \\ \hline
\end{tabularx}

A scrambled middle does not disqualify a signal: clean tails around a muddled G2--G4 is a common
shape and a perfectly tradeable one, since the tails are where the positions go.

On direction, adopt a convention and keep it: \textbf{leave the signal's sign alone and let the
chart come out upside down.} Flipping a signal the moment its staircase descends is cheap, and a
dozen signals later you will not remember which ones you flipped. Read a descending staircase as
\emph{this works, traded short}, and carry the sign in the strategy rather than in the signal
definition.

\begin{claim}\label{clm:averaging}
The ordering surfaces at all only because of step 4---averaging shrinks the noise and leaves the
signal untouched.
\end{claim}

\begin{proof}
Take $m$ observations sharing a similar signal value. Their mean return still has expectation
$\beta$ times their mean signal---they were chosen for having nearly the same signal, so averaging
them changes it barely at all---while the noise around it falls to $\sigma_\epsilon / m^{1/2}$.
Measured, with the edge taken as the G5 $-$ G1 gap the chart is trying to resolve:

\begin{center}
\begin{tabular}{|l|c|c|}
\hline
       & \textbf{One observation} & \textbf{Mean of 11,522} \\ \hline
Signal & 10.4 bp & 10.4 bp \\ \hline
Noise  & 149 bp  & $149 / 11522^{1/2} \approx 1.4$ bp \\ \hline
Ratio  & 1 : 14  & $\mathbf{7 : 1}$ \\ \hline
\end{tabular}
\end{center}

The left column is the scatter of \secref{sec:hoped-vs-got} and the right column is one bar of the
chart. Nothing was added---only the noise was taken away.
\end{proof}

\fig{noise-shrinks}{The same measured population as four bucket charts, mean forward return in basis
points against signal bucket, drawn from 30, 300, 3,000 and all 11,522 observations per bucket. At
$m = 30$ the bars swing over tens of basis points with error bars wider still and no ordering; by
$m = 3{,}000$ G1 is visibly the tallest bar; at $m = 11{,}522$ the error bars are a few basis points
and G1 stands clear of its own error while G2 to G5 sit near zero.}{fig:noise-shrinks}

The quantity being estimated is identical in all four panels---the real bucket means, $+11$ bp at G1
and roughly zero at G5. Only the error moves, and at $m = 30$ it is five times the whole staircase.
\textbf{Sample size is not a detail of the recipe; it is the reason the recipe works.}

\begin{note}[Where that overstates it]
The $m^{1/2}$ assumes independence, and one asset's history violates it twice over: overlapping
lookback windows share most of their inputs, and returns cluster in time, so consecutive
observations are near-copies rather than independent draws. Both push the effective count well below
$m$, and the noise does not really reach 1.4 bp. The logic survives: averaging cancels noise and
leaves systematic signal.
\end{note}

\subsubsection{Is the staircase real?}\label{sec:staircase-real}

\secref{sec:bar-plot-trend} read the rise against the error bars by eye. That comparison is a
hypothesis test, and it is worth writing down, because the eyeball version and the arithmetic
version do not agree.

\begin{definition}[Standard error of a bucket]\label{def:se-bucket}
For bucket $g$ holding $m_g$ observations whose forward returns have sample standard deviation
$\sigma_g$,
\[
  \mathrm{SE}_g = \frac{\sigma_g}{m_g^{1/2}}
\]
the spread of that bucket's \textbf{mean}, not of the observations inside it---which is why it
shrinks with $m_g$ while $\sigma_g$ does not. Bars here are drawn at $\pm 2\,\mathrm{SE}_g$.
\end{definition}

\textbf{The hypothesis on trial is \secref{sec:signal-hypothesis}'s, narrowed to the two buckets
that carry the positions.} A long book buys G5 and sells G1, so what has to be non-zero is the
difference between those two means:
\[
  H_0: \mu_{G5} = \mu_{G1}
\]
against the alternative that they differ, where $\mu_g$ is bucket $g$'s true mean forward return.
The statistic is the observed difference in units of its own noise:
\[
  z = \frac{\mu_{G5} - \mu_{G1}}{\left( \mathrm{SE}_{G5}^2 + \mathrm{SE}_{G1}^2 \right)^{1/2}}
\]
the denominator adding \textbf{variances} rather than standard errors, because the two bucket means
are computed from disjoint sets of dates.

\begin{note}[Why it is a $t$ and gets called a $z$]
Both $\mu_g$ and $\sigma_g$ are estimated from the same sample, so the exact null distribution is
Student's $t$, not the normal. It stops mattering fast: at bucket sizes in the thousands the two are
indistinguishable and the number is quoted as a \textbf{z-score}. Derive it as a $t$, report it as a
$z$, and at $m_g$ in the tens use the $t$ and mean it.
\end{note}

\begin{claim}\label{clm:z-two}
$\left| z \right| \geq 2$ rejects $H_0$ at the 5 percent level.
\end{claim}

\begin{proof}
Under $H_0$ the statistic is standard normal in the large-sample limit, and a standard normal places
0.0455 of its mass outside $\pm 2$---less than $\alpha = 0.05$, so the observation lies inside the
two-sided rejection region.
\end{proof}

\begin{example}\label{ex:measured-z}
The whole sample, five buckets, 11,522 dates in the smallest:

\begin{center}
\begin{tabular}{|c|r|r|r|r|}
\hline
\textbf{Bucket} & $\mu_g$ & $\sigma_g$ & $m_g$ & $\mathrm{SE}_g$ \\ \hline
G1 & $\mathbf{+10.72}$ \textbf{bp} & 173 bp & 11,522 & 1.62 bp \\ \hline
G2 & $+1.05$ bp & 144 bp & 11,523 & 1.34 bp \\ \hline
G3 & $+3.48$ bp & 136 bp & 11,525 & 1.27 bp \\ \hline
G4 & $+2.83$ bp & 140 bp & 11,523 & 1.30 bp \\ \hline
G5 & $\mathbf{+0.31}$ \textbf{bp} & 149 bp & 11,556 & 1.39 bp \\ \hline
\end{tabular}
\end{center}

\[
  z = \frac{0.31 - 10.72}{\left( 1.39^2 + 1.62^2 \right)^{1/2}} = \frac{-10.41}{2.13} \approx -4.9
\]

Well past 2 in absolute value, and the bars---G1 spanning $+7.5$ to $+14.0$ bp, G5 spanning $-2.5$
to $+3.1$---clear each other with room to spare.
\end{example}

\begin{note}[Non-overlap is sufficient, not necessary]
The two readings differ, and in the direction that costs you signals. Bars of half-width
$2\,\mathrm{SE}$ fail to touch when the difference exceeds
$2 \left( \mathrm{SE}_{G5} + \mathrm{SE}_{G1} \right)$, while the test rejects when it exceeds
$2 \left( \mathrm{SE}_{G5}^2 + \mathrm{SE}_{G1}^2 \right)^{1/2}$---and a sum is never smaller than
the root of the sum of squares. On the measured numbers the eye demands 6.0 bp and the test demands
4.3 bp, so a difference between those two is one the picture calls inconclusive and the arithmetic
calls significant. \textbf{Bars that clear each other are always significant; bars that touch may
still be.} Use the chart to accept and compute $z$ before rejecting.
\end{note}

\textbf{The sign is negative, and that is a result rather than a failure.} $z = -4.9$ says the two
means differ, decisively; it says nothing about which way round they should be. What was measured is
that \textbf{a high 21-day momentum is followed by a slightly worse next session than a low one}
---the descending staircase \secref{sec:bar-plot-trend} told you to expect and to keep. It is not a
broken signal, it is \textbf{reversal}, and the Background section explains where it comes from and
what removes it: put a gap between the signal's window and the return it is scored on, and the sign
comes back.

\begin{note}[What the test does not license]
Two things sit outside it. The $m_g^{1/2}$ in every standard error assumes independent observations,
and \secref{sec:bar-plot-trend}'s closing note explains why one asset's overlapping windows are
nothing of the kind---the effective count is well below $m_g$, so every $z$ here is optimistic. And
a threshold of 5 percent means one signal in twenty passes on noise, which is a fact about
\emph{one} test: running twenty variants and reporting the one that cleared 2 is not a test at all,
and is \chapref{5 \midot{} Modelling}{05_modelling}'s subject.
\end{note}

\textbf{Failing to reject is a verdict on the measurement, not on the signal.}
$\left| z \right| < 2$ says the data did not separate the two means. Whether that is because they
are equal, or because the error bars were too wide to tell them apart, is a question about $m_g$
---and the two readings are not interchangeable. The smallest gap the test could have caught, at the
four sample sizes of the figure above and a common $\sigma_g$ of 149 bp:

\begin{tabularx}{\linewidth}{|r|>{\raggedright\arraybackslash}p{0.11\linewidth}|>{\raggedright\arraybackslash}p{0.17\linewidth}|X|}
\hline
$m_g$ & $\mathrm{SE}_g$ & \textbf{Smallest detectable G5 $-$ G1} & \textbf{What a null result there rules out} \\ \hline
30     & 27 bp   & 77 bp  & Almost nothing. An edge large enough to build a career on is still
                            consistent with it \\ \hline
300    & 8.6 bp  & 24 bp  & Only edges several times larger than anything this chapter
                            measures \\ \hline
3,000  & 2.7 bp  & 7.7 bp & Most of the range that matters, but not the part costs decide \\ \hline
11,522 & 1.4 bp  & 3.9 bp & The effect, down to a gap too small to survive execution anyway \\ \hline
\end{tabularx}

Only the bottom row turns a null result into a finding. In the top row the gate has not been failed
so much as attempted---\textbf{the honest report is ``not measured yet'', and the fix is more
observations, not a different signal.}

\subsubsection{What the bar plot cannot show}\label{sec:bar-plot-cannot}

Every date contributes one observation---its signal's score, and the return that followed. Step 5
flattens all of them onto one picture, five slots wide.

\fig{bucket-time-collapse}{Left, one asset's dates strung along a time axis, each drawn at the slot
its signal scored into, so all five rows are populated throughout and a violent stretch at the right
is shaded. Right, the same observations with the date discarded: five bars of mean forward return
with error bars. The arrow between them averages down each column, and the time axis is what it
spends.}{fig:bucket-time-collapse}

Step 5 integrates over $t$, and an integral hands back an area, never the shape of the function
underneath it. A staircase that is monotone over twenty years is equally consistent with an edge
that held throughout and with one that worked for five years and was flat for fifteen.

\begin{tabularx}{\linewidth}{|>{\raggedright\arraybackslash}p{0.34\linewidth}|X|}
\hline
\textbf{Not on the chart} & \textbf{Where to look instead} \\ \hline
\textbf{When} the edge happened
  & the equity curve of \chapref{5 \midot{} Modelling}{05_modelling} \\ \hline
\textbf{Who} is inside a bucket---which dates, which regime
  & the dates behind each bar, tabulated rather than averaged \\ \hline
\textbf{The spread} of returns behind a bar, as against its mean
  & the distribution within a bucket, not its mean \\ \hline
\textbf{How independent} the observations are
  & overlapping windows, per the note above \\ \hline
\end{tabularx}


\section{Step 2, corrected --- risk-adjusted momentum}\label{sec:risk-adjusted}

\textbf{This section repairs \secref{sec:beta-method}, it does not extend it.} Every piece of
machinery there survives: the sort key is knowable at $t$, the returns are recorded honestly, and
the averaging shrinks noise exactly as the proof says. What \secref{sec:beta-method} never settled
is \emph{what step 2 ranks}, and ranking the raw momentum value produces a chart about something
other than the signal.

Raw momentum is not comparable \textbf{across time}, and \secref{sec:beta-method} compares nothing
else: its five slots sort one asset's own dates against one another. 2 percent in the calm 2021 tape
and 2 percent in the 2023 rate-hike drawdown are different events, and filing them into the same
slot says they are the same.

The identical failure reappears one step later, \textbf{across assets}, the moment you want more
than one name in a book---and there it is an assumption you have to state rather than a fact. A
universe of five hundred large-cap equities might plausibly sit on one scale already. A CTA universe
does not: 50 bp in a session is an ordinary day for an equity index and, for a Treasury future, an
event that would lead the news. 2 percent monthly momentum in a Treasury ETF is a large move; the
same 2 percent in a semiconductor ETF is a quiet month. Rank the two against each other on raw
momentum and the semiconductor wins every time---not because its trend is stronger but because
everything about it is larger.

Divide by volatility:
\[
  \mathit{MOM}^{\text{risk-adj}}_{s,t}
    = \operatorname{Avg}\left( \frac{r_{s,t-i}}{\sigma_{s,t}} \right), \qquad i = 1 \ldots N
\]
where $\sigma_{s,t}$ is that asset's volatility, estimated on data strictly before $t$---a
denominator that peeks at the future contaminates the signal as surely as a numerator would
(\secref{sec:info-availability}).

Every date, and later every asset, now lands on one scale, so \secref{sec:beta-method}'s slots
compare like with like---two students both scoring 80 on different exams against different cohorts
have not achieved the same thing.

\begin{note}[A second route to the same plane]
Dividing by $\sigma_{s,t}$ is not the only way. You can also replace the value outright by \textbf{its
percentile against that asset's own past}, which lands every date on $[0, 1]$ by construction and
needs no volatility estimate at all. Either route hands step 2 something it can legitimately rank.
\end{note}

\begin{example}\label{ex:percentile}
Two assets, each with five past signal values and one for today:

\begin{center}
\begin{tabular}{|c|l|c|c|c|}
\hline
\textbf{Asset} & \textbf{Its own past signals} & \textbf{Today} & \textbf{Beats} & \textbf{Percentile} \\ \hline
A & $-10$, 8, 4, 0, $-5$          & $\mathbf{+7}$  & four of five & \textbf{80} \\ \hline
B & $-100$, $-80$, 90, $-60$, 20  & $\mathbf{-70}$ & two of five  & \textbf{40} \\ \hline
\end{tabular}
\end{center}

Raw, $+7$ against $-70$ is not a comparison anybody should make. As percentiles, 80 against 40
is---and it says A is the stronger trend \emph{relative to its own history}, which is the only sense
in which the question has an answer. The cost is the one every rank charges: the magnitudes are
gone.
\end{example}

\begin{note}[Which cut to use once it is scored]
The percentile route hands back a uniform score, so cutting at the quintiles gives five equal groups
for free. The volatility route does not: a standardized signal is roughly normal, so cutting it at
fixed intervals such as $\pm 1$ and $\pm 2$ leaves almost everything in the middle and starves the
two tails you would actually trade. Cut it at its quintiles instead.
\end{note}

\textbf{The rule this section leaves behind.} Every signal is \textbf{standardized} or \textbf{scored
against its own past} before it is compared to anything---otherwise the bar plot reports volatility
while wearing the signal's name. The measured chart back in \secref{sec:beta-method} already
followed this rule before this section stated it: the percentile route is what let 37 assets share
one pool of buckets.


\section{What to report --- and why it is not a Sharpe ratio}\label{sec:what-to-report}

Every chart in \secref{sec:two-ways-out} reports one quantity: \textbf{the mean forward return of a
bucket, with its standard error.} The instinct at this point is to reach for something that sounds
more like a result, and the one everybody reaches for is the wrong tool for this stage.

\begin{definition}[Sharpe ratio]\label{def:sharpe}
A portfolio's average return in excess of what financing it costs, per unit of the volatility it
ran:
\[
  \mathrm{SR} = \frac{E[R_p] - r_f}{\sigma_p}
\]
with $R_p$ the portfolio's return over a period, $r_f$ the risk-free rate charged on the capital it
used, and $\sigma_p$ the standard deviation of $R_p$.
\end{definition}

It is the right measurement for a finished portfolio and the wrong one here, for four separate
reasons.

\textbf{There is no portfolio yet.} \secref{sec:signal-is-portfolio}'s argument is precisely that
the signal can be judged \emph{without} building one, and nothing in \secref{sec:two-ways-out} has
chosen a weight. A Sharpe ratio scores the output of a construction step that has not happened;
computing one on a bucket's returns scores a portfolio nobody would run---equal weight on every date
that fell into one fifth of the sample.

\textbf{The financing rate is not zero here, and it is not constant.} The numerator subtracts $r_f$,
and which rate that is depends on the book's \textbf{net exposure}
(\secref{sec:signal-is-portfolio}):

\begin{tabularx}{\linewidth}{|>{\raggedright\arraybackslash}p{0.30\linewidth}|>{\raggedright\arraybackslash}p{0.14\linewidth}|X|}
\hline
\textbf{Book} & \textbf{Net exposure} & \textbf{Correct $r_f$} \\ \hline
Dollar-neutral long/short---one dollar long against one dollar short
  & $\approx 0$
  & \textbf{zero.} The short leg's proceeds fund the long leg, so no capital is being
    borrowed \\ \hline
A 150/50 CTA book
  & 100 percent
  & \textbf{the actual financing rate}, on the capital the net exposure represents \\ \hline
\end{tabularx}

Statistical-arbitrage books are the first row, which is where the habit of setting $r_f = 0$ comes
from and where it is correct. A CTA book is the second row, and the rate belonging in it is a
\textbf{time series}, not a constant: the policy rate this sample spans moved from near zero to
above five percent. Carrying over ``we always used 2 percent'' from a book that was dollar-neutral
puts an invented series into the numerator of every number reported.

\textbf{It folds three estimates into one number.} Mean, volatility and financing rate, each with
its own error. A Sharpe that falls because the edge died and a Sharpe that falls because the tape
got noisier are the same number---the exact opposite of \secref{sec:beta-method}'s method, where one
chart carries one quantity so a change has one cause.

\textbf{A Sharpe means nothing without the market it was earned in.} Two books, and the smaller
number is the better strategy:

\begin{tabularx}{\linewidth}{|>{\raggedright\arraybackslash}p{0.17\linewidth}|>{\raggedright\arraybackslash}p{0.13\linewidth}|>{\raggedright\arraybackslash}p{0.13\linewidth}|X|}
\hline
\textbf{The tape} & \textbf{Market's own Sharpe} & \textbf{The book's Sharpe} & \textbf{Reading} \\ \hline
Violent throughout & 0.4 & \textbf{0.8}
  & held together through amplitude that should have wrecked it, at twice the market's own
    risk-adjusted return---an excellent result \\ \hline
Calm and rising & 0.8 & \textbf{0.9}
  & barely beat holding the index. The extra 0.1 is within what one lucky bet would
    produce \\ \hline
\end{tabularx}

\textbf{The mean return is the quantity you can actually trade.} It is what the book delivers, it is
what \secref{sec:signal-is-portfolio}'s algebra is written in, and it is denominated in the same
basis points as the costs that will be subtracted from it later. A Sharpe ratio is a score attached
afterwards, useful for comparing \textbf{finished} portfolios against each other---which belongs to
performance evaluation, after an optimization step has given it something to measure.


\section{Information availability}\label{sec:info-availability}

Everything above assumes the signal at $t$ uses only data knowable at $t$. Look-ahead bias is born
here; the execution offsets of \chapref{5 \midot{} Modelling}{05_modelling} are the second line of
defense and cannot rescue a signal contaminated at construction. The rolling-quantile rule
(\secref{sec:risk-adjusted}) and the train/validation/test split
(\chapref{5 \midot{} \S 1}{05_modelling}) are the same discipline.


\section*{Background}
\addcontentsline{toc}{section}{Background}

\subsection*{Why does a trend hand part of itself back, and can that be competed away?}

\textbf{Reversal is not an anomaly, it is the normal state of the tape.} A trend that has just
formed hands a little of it back, and at daily frequency that rebate is large enough to cancel the
edge outright---which is why \secref{sec:which-return} leaves a gap $g$ between the signal's window
and the return it is scored on. Three mechanisms produce it, and they are not equally fragile:

\begin{tabularx}{\linewidth}{|>{\raggedright\arraybackslash}p{0.20\linewidth}|X|>{\raggedright\arraybackslash}p{0.22\linewidth}|}
\hline
\textbf{Mechanism} & \textbf{What happens} & \textbf{Competed away?} \\ \hline
\textbf{Bid--ask bounce}
  & Closes print alternately at the bid and the ask, oscillating around the true mid---negative
    serial correlation by construction
  & \textbf{No.} A measurement artefact; no price moved \\ \hline
\textbf{Paying for liquidity}
  & An order that must fill now pushes price past fair value, and whoever takes the other side is
    repaid when it returns
  & \textbf{No.} A fee for a service actually rendered \\ \hline
\textbf{Overreaction}
  & News is over-extrapolated, then partly corrected
  & \textbf{Yes}---the only part that decays \\ \hline
\end{tabularx}

Two of the three are not mistakes, which is why reversal survives competition instead of being
arbitraged into nothing---only the overreaction row decays as more people trade against it.


\section*{Appendix \midot{} Notation}
\addcontentsline{toc}{section}{Appendix \midot{} Notation}

Throughout, $s$ indexes the asset and $t$ the date. The signal is computed per asset, so $s$ is
dropped wherever a formula never leaves one asset; it carries weight everywhere the assets are
ranked against each other.

\begin{xltabular}{\linewidth}{|>{\raggedright\arraybackslash}p{0.22\linewidth}|X|>{\raggedright\arraybackslash}p{0.16\linewidth}|}
\hline
\textbf{Symbol} & \textbf{Means} & \textbf{First used} \\ \hline
\endhead
$s$, $t$
  & the asset, and the date in periods (days here)
  & \secref{sec:signal-hypothesis} \\ \hline
$r_{s,t}$, $\mathit{MOM}_{s,t}$, $w_{s,t}$
  & that asset's return in that period, the momentum signal it produces, and the share of capital
    that signal earns it
  & \secref{sec:three-objects} \\ \hline
$R_t$
  & the portfolio's return on that date, $\sum_s w_{s,t} r_{s,t}$
  & \secref{sec:signal-is-portfolio} \\ \hline
$G$, $\delta$
  & the gross-exposure target the weights must sum in absolute value to, and the tolerance allowed
    on the net
  & \secref{sec:signal-is-portfolio} \\ \hline
$\rho_s$, $\sigma_{\mathit{MOM},s}$, $\sigma_{r,s}$
  & one asset's signal-return correlation, and the standard deviations of its signal and its
    forward return
  & \secref{sec:signal-is-portfolio} \\ \hline
$N$, $i$
  & lookback length, and the lag inside it running $1$ to $N$
  & \secref{sec:three-objects} \\ \hline
$x$, $y$, $\beta$, $\epsilon$
  & one observation's signal value and forward return, the slope between them, and the part of the
    return the signal cannot reach
  & \secref{sec:why-cloud} \\ \hline
$\rho$, $R^2$
  & their correlation, and its square---the share of the return's variance the signal explains
  & \secref{sec:why-cloud} \\ \hline
$\sigma_x$, $\sigma_y$, $\sigma_\epsilon$
  & standard deviation of the signal, of the return, and of the unreachable part
  & \secref{sec:why-cloud} \\ \hline
$m$, $m_g$
  & observations sharing a bucket, and the count in bucket $g$
  & \secref{sec:beta-method} \\ \hline
$\mu_g$, $\sigma_g$, $\mathrm{SE}_g$
  & one bucket's true mean forward return, the sample standard deviation of the returns in it, and
    the standard error of its mean
  & \secref{sec:staircase-real} \\ \hline
$H_0$, $z$, $\alpha$
  & the null that two buckets share a mean, the standardized G5 $-$ G1 difference testing it, and
    the level it is tested at
  & \secref{sec:staircase-real} \\ \hline
G1 \ldots G5
  & the buckets, lowest to highest signal
  & \secref{sec:beta-method} \\ \hline
$\sigma_{s,t}$
  & one asset's volatility, estimated on data before that date
  & \secref{sec:risk-adjusted} \\ \hline
$R_p$, $r_f$, $\sigma_p$
  & a portfolio's return over a period, the financing rate subtracted from it, and its return
    volatility
  & \secref{sec:what-to-report} \\ \hline
$g$
  & gap length---periods between the end of the signal's window and the return it is scored on
  & \secref{sec:which-return} \\ \hline
\end{xltabular}

\begin{note}[Collisions to watch]
$\alpha$ is the test level here, and \texttt{alpha} in code font is matplotlib's opacity
(\secref{sec:alpha-opacity}); neither is \chapref{3}{03_shaping_the_lookback}'s smoothing constant.
$\sigma$ carries whatever its subscript says: $\sigma_g$ is the spread of returns \emph{inside} a
bucket, $\mathrm{SE}_g$ the spread of that bucket's mean, and $\sigma_{s,t}$ an asset's trailing
volatility---not \chapref{5 \midot{} \S 4.2}{05_modelling}'s $\sigma_j$, which rescales one signal.
$\beta$ is the slope of return on signal, never a market beta. $g$ is the gap in
\secref{sec:which-return} and a bucket index in \secref{sec:staircase-real}; the subscript position
tells them apart.
\end{note}
